\chapter{Introduction}

\section{Background, Motivation and Scope}
Machine learning turns large, noisy and often unlabelled data into decisions. Most real problems call for more than one
kind of technique. First the data has to be understood and cleaned. Then predictive models are built and compared.
Hidden structure has to be discovered where no labels exist, and time-ordered behaviour has to be modelled. Finally,
reliable high-performing models are built by combining many learners. This project applies the complete toolkit of the
21CSC305P Machine Learning course across three Project-Based Learning (PBL) reviews and five units. Each unit works on a
dataset chosen for the problem it addresses.

The application domains are deliberately varied:
\begin{itemize}
\item \textbf{The digital economy (PBL I):} mobile-app and video-game markets, and the direct-marketing campaigns of a bank.
App developers, publishers and marketing teams need to know what drives engagement, sales and customer response.
\item \textbf{Energy (PBL I):} accurate electrical-output prediction for a power plant supports grid operators in planning
dispatch.
\item \textbf{Ubiquitous health sensing (PBL II):} recognising a person's physical activity from a smartphone accelerometer
is the basis of fitness tracking, elderly-care monitoring and rehabilitation.
\item \textbf{Clinical decision support (PBL III):} early identification of patients at risk of diabetes from routine
health records lets clinicians intervene before complications develop. Diabetes affects hundreds of millions of people
worldwide \cite{ada,miotto}.
\end{itemize}
The main stakeholders are product and marketing teams, energy operators, health-technology developers, clinicians and
patients. The scope covers data acquisition, cleaning, exploratory analysis, model implementation, rigorous evaluation on
held-out data and interpretation of the results. Deployment is outside the scope and is discussed as future work.

\section{Integrated Problem Statement}
The overall problem is to \emph{design, implement and evaluate machine-learning pipelines that turn raw data into reliable,
interpretable decisions}, and to show how the choice of technique depends on the structure of the data. The problem is
broken into five linked sub-problems, one per unit:
\begin{enumerate}
\item \textbf{Unit 1 (Exploratory analysis):} what drives success in the mobile-app and video-game markets?
\item \textbf{Unit 2 (Supervised learning):} predict a continuous quantity (power-plant energy output) and a binary
business outcome (whether a bank client subscribes to a term deposit).
\item \textbf{Unit 3 (Unsupervised learning):} discover groups of similar sensor windows that correspond to physical
activities, without using labels.
\item \textbf{Unit 4 (Sequential modelling):} infer the hidden sequence of a person's activities from accelerometer time
series using Markov models, Hidden Markov Models, forward–backward and Viterbi.
\item \textbf{Unit 5 (Ensemble learning):} build an intelligent diabetes prediction model from patient health records and
investigate how effective different ensemble classifiers are compared with single-model baselines.
\end{enumerate}

\section{Project Objectives}
\begin{itemize}
\item \textbf{Objective 1 (Unit 1, EDA):} clean and explore two real datasets (10,841 Play Store apps and 16,598 video
games) and extract quantified, actionable insights.
\item \textbf{Objective 2 (Unit 2, regression and classification):} compare at least ten regressors and ten classifiers
with cross-validation. Expected result: a tuned model clearly better than the linear baseline, and an imbalance-aware
classifier that avoids the accuracy trap.
\item \textbf{Objective 3 (Unit 3, clustering):} apply K-Means and Gaussian Mixture clustering, with PCA for
visualisation, to accelerometer features, and measure how well the clusters recover the true activities.
\item \textbf{Objective 4 (Unit 4, sequential modelling):} implement forward–backward and Viterbi from first principles,
verify them against a library, and decode activity sequences of unseen people with unsupervised and supervised HMMs.
\item \textbf{Objective 5 (Unit 5, ensemble learning):} implement at least three ensemble algorithms (ten are
implemented) and compare them with Logistic Regression and Decision Tree baselines using F1-score and ROC-AUC. Expected
result: an ensemble model with high F1 and ROC-AUC on unseen patients.
\end{itemize}

\section{Unit-wise Topic and Dataset Allocation}
Table~\ref{tab:allocation} lists the topic and dataset of every unit. The datasets were allotted per review, so each
chapter documents its own dataset, preprocessing and EDA.

\begin{table}[H]
\centering
\caption{Unit-wise topic and dataset allocation}
\label{tab:allocation}
\renewcommand{\arraystretch}{1.3}
\begin{tabularx}{\linewidth}{@{}l c >{\raggedright\arraybackslash}X >{\raggedright\arraybackslash}X@{}}
\toprule
\textbf{PBL} & \textbf{Unit} & \textbf{Allotted topic / algorithms} & \textbf{Dataset} \\
\midrule
PBL 1 & 1 & Exploratory Data Analysis & Google Play Store Apps; Video Game Sales (Kaggle) \\
PBL 1 & 2 & Regression and classification & Combined Cycle Power Plant; Bank Marketing (UCI) \\
PBL 2 & 3 & Clustering: K-Means, Gaussian Mixture, PCA & UCI HAR (window features) \\
PBL 2 & 4 & Markov models, HMM, forward–backward, Viterbi & UCI HAR (accelerometer time series) \\
PBL 3 & 5 & Ensemble learning: bagging, boosting, voting, stacking & Diabetes Prediction Dataset (Kaggle) \\
\bottomrule
\end{tabularx}
\end{table}

The rest of the report is organised as follows. Chapter 2 covers Units 1 and 2, Chapter 3 Units 3 and 4, and Chapter 4
Unit 5. Chapter 5 brings the findings together, and Chapter 6 concludes. Appendix A contains the complete code.
